% !TeX root=main.tex
% دستور زیر باعث شماره‌بندی صفحات فصول از ۱ می‌شود و باید در اولین فصل شما باشد. آن را حذف نکنید!
\pagenumbering{arabic} % 1, 2, ...

\chapter{مقدمه}
\section{مقدمه کلی}

بلاکچین\LTRfootnote{Blockchain} را می‌توان یکی از مهم‌ترین نوآوری‌های ترکیبی در حوزه سامانه‌های توزیع‌شده، رمزنگاری کاربردی و اقتصاد شبکه دانست. این فناوری با کنار هم قرار دادن چند ایده شناخته‌شده مانند امضای دیجیتال، تابع هش رمزنگاری، زنجیره بلوکی، شبکه همتا‌به‌همتا، انگیزه اقتصادی و اجماع توزیع‌شده، امکان نگهداری یک دفترکل مشترک بدون متولی مرکزی را فراهم می‌کند. در چنین سامانه‌ای، کاربران برای انتقال ارزش یا ثبت وضعیت نیازی ندارند به بانک، شرکت پرداخت، مرکز داده یا سرور قابل اعتماد واحدی تکیه کنند؛ بلکه اعتبار وضعیت سامانه از قواعد رمزنگاری و توافق جمعی نودها حاصل می‌شود.

بیتکوین\LTRfootnote{Bitcoin} نخستین نمونه موفق و پایدار این رویکرد است. در طراحی بیتکوین، تراکنش‌ها به صورت عمومی اعلام می‌شوند، نودها اعتبار آن‌ها را بررسی می‌کنند، ماینرها تراکنش‌های معتبر را در بلوک قرار می‌دهند، و با حل یک مسئله اثبات کار\LTRfootnote{Proof of Work (PoW)}، بلوک پیشنهادی خود را به شبکه ارائه می‌کنند. زنجیره‌ای که بیشترین کار انباشته را در خود دارد به عنوان تاریخچه معتبرتر پذیرفته می‌شود. این مدل باعث می‌شود تغییر یک تراکنش قدیمی بسیار پرهزینه باشد، زیرا مهاجم باید کار محاسباتی بلوک هدف و تمام بلوک‌های بعد از آن را دوباره انجام دهد و از مجموع توان محاسباتی نودهای صادق پیشی بگیرد.

با وجود موفقیت امنیتی و اقتصادی بیتکوین، معماری آن از دید مقیاس‌پذیری محدود است. در یک زنجیره عمومی کاملاً تکرارشونده، همه نودهای کامل باید تمام تراکنش‌ها را دریافت، ذخیره و اعتبارسنجی کنند. این ویژگی، اگرچه برای امنیت و استقلال نودها ضروری است، نرخ پردازش تراکنش را محدود می‌کند. به بیان ساده، اگر ظرفیت هر بلوک محدود باشد و بلوک‌ها با فاصله زمانی نسبتاً زیاد تولید شوند، تعداد تراکنش‌های قابل پردازش در واحد زمان نیز محدود خواهد بود. این مسئله به ویژه در زمان افزایش تقاضا خود را نشان می‌دهد: تراکنش‌ها در استخر تراکنش\LTRfootnote{Mempool} انباشته می‌شوند، کاربران برای ورود سریع‌تر به بلوک کارمزد بیشتری پیشنهاد می‌کنند، و کاربردهایی مانند پرداخت خرد، بازی‌های بلاکچینی، اینترنت اشیا و بازارهای مالی پرسرعت با مشکل مواجه می‌شوند.

در نتیجه، پرسش اصلی این پایان‌نامه چنین است: چگونه می‌توان توان عملیاتی\LTRfootnote{Throughput} بلاکچین را افزایش داد، بدون آنکه امنیت زنجیره اصلی و اصل غیرمتمرکز بودن به طور جدی قربانی شود؟ پاسخ پیشنهادی این پژوهش، استفاده از معماری چندزنجیره‌ای است که در آن زنجیره اصلی نقش امنیتی و هماهنگ‌کننده دارد و چندین زنجیره جانبی\LTRfootnote{Sidechain} مستقل، پردازش تراکنش‌ها را موازی می‌کنند.

\section{بیان مسئله}

مسئله مقیاس‌پذیری در بلاکچین فقط یک مسئله مربوط به افزایش سرعت برنامه‌نویسی یا ارتقای سخت‌افزار نیست؛ بلکه مسئله‌ای سیستمی و چندبعدی است. در سامانه‌های متمرکز، افزایش ظرفیت معمولاً با اضافه کردن سرور، استفاده از پایگاه‌داده توزیع‌شده، افزایش پهنای باند یا بهینه‌سازی نرم‌افزار انجام می‌شود. اما در بلاکچین‌های عمومی، هر تغییر در ظرفیت باید با ملاحظات امنیتی و تمرکززدایی همراه باشد. اگر اندازه بلوک بسیار بزرگ شود، فقط نودهایی با منابع سخت‌افزاری و شبکه‌ای قوی می‌توانند کل زنجیره را نگهداری کنند. اگر فاصله بلوک بسیار کوتاه شود، احتمال انتشار دیرهنگام بلوک‌ها، ایجاد فورک‌های موقت\LTRfootnote{Fork} و افزایش نرخ بلوک‌های کهنه\LTRfootnote{Stale Block} بالا می‌رود. اگر اجماع به تعداد کمی اعتبارسنج واگذار شود، توان عملیاتی افزایش می‌یابد، اما شبکه به سمت اعتماد به گروه کوچک‌تری از بازیگران حرکت می‌کند.

برای آنکه مسئله به شکل دقیق قابل تحلیل باشد، ابتدا باید مدل سامانه را رسمی کنیم. بلاکچین را می‌توان به صورت یک ماشین حالت تکرارشونده\LTRfootnote{Replicated State Machine (RSM)} توصیف کرد: مجموعه‌ای از نودها که هر یک نسخه‌ای از وضعیت \(\sigma\) را نگه می‌دارند و پروتکل اجماع، ترتیب کلی یکتایی روی تراکنش‌ها تعیین می‌کند، به گونه‌ای که اعمال دنباله مرتب‌شده تراکنش‌ها روی وضعیت اولیه، در همه نودهای صادق به وضعیت یکسانی برسد:

\begin{latin}
\[
\sigma_{t+1} = \Gamma(\sigma_t, tx_t), \qquad \forall\, n_i, n_j \in H:\; \sigma^{(i)} = \sigma^{(j)}
\]
\end{latin}

که در آن \(\Gamma\) تابع گذار وضعیت (اعتبارسنجی و اعمال تراکنش)، و \(H\) مجموعه نودهای صادق است. در بیتکوین، وضعیت \(\sigma\) همان مجموعه خروجی‌های تراکنش خرج‌نشده\LTRfootnote{Unspent Transaction Output (UTXO)} است و تابع \(\Gamma\) شامل بررسی امضا، وجود ورودی‌ها در مجموعه خروجی‌های خرج‌نشده و برقراری موازنه مقادیر است. نکته کلیدی این است که در معماری تک‌زنجیره‌ای، همه نودهای کامل باید کل دنباله تراکنش‌ها را اجرا کنند؛ یعنی ظرفیت کل سامانه برابر ظرفیت کندترین نودِ قابل قبول است و افزودن نود جدید، برخلاف سامانه‌های توزیع‌شده متعارف، ظرفیت را افزایش نمی‌دهد، فقط افزونگی را بالا می‌برد.

با این مدل، می‌توان گلوگاه‌های فنی را به تفکیک شناسایی کرد. گلوگاه نخست، پهنای باند و تأخیر انتشار است: بلوکی به اندازه \(B\) بایت باید در بازه‌ای کوتاه به اکثریت نودها برسد. اگر \(\Delta\) تأخیر مؤثر انتشار بلوک در شبکه و \(T\) فاصله متوسط تولید بلوک باشد، احتمال آنکه در بازه انتشار یک بلوک، بلوک رقیبی تولید شود (و فورک موقت رخ دهد) با تقریب فرایند پواسون چنین است:

\begin{latin}
\[
p_{fork} \approx 1 - e^{-\Delta / T}
\]
\end{latin}

این رابطه نشان می‌دهد چرا نمی‌توان \(T\) را آزادانه کوچک یا \(B\) را (که \(\Delta\) را بزرگ می‌کند) آزادانه بزرگ کرد: در هر دو حالت \(\Delta/T\) رشد می‌کند، نرخ فورک بالا می‌رود، بخشی از توان هش صرف بلوک‌هایی می‌شود که به زنجیره نهایی نمی‌رسند، و امنیت مؤثر در برابر مهاجم کاهش می‌یابد. گلوگاه دوم، هزینه اعتبارسنجی است: بررسی هر امضای دیجیتال منحنی بیضوی\LTRfootnote{Elliptic Curve Digital Signature Algorithm (ECDSA)} چند ده میکروثانیه زمان پردازنده می‌گیرد و تراکنشی با چند ورودی، چند بار این هزینه را تحمیل می‌کند؛ در نرخ‌های بالای تراکنش، همین هزینه به تنهایی پردازنده نودهای معمولی را اشباع می‌کند. گلوگاه سوم، دسترسی تصادفی به وضعیت است: مجموعه خروجی‌های خرج‌نشده بیتکوین چند گیگابایت است و الگوی دسترسی به آن تصادفی است؛ بنابراین با رشد نرخ تراکنش، ورودی/خروجی دیسک و حافظه نهان به گلوگاه تبدیل می‌شوند. هر راهکار مقیاس‌پذیری باید مشخص کند کدام‌یک از این سه گلوگاه را برطرف می‌کند و کدام را فقط جابه‌جا می‌کند.

در بیتکوین، ظرفیت لایه پایه تابعی از سه عامل کلیدی است: اندازه یا وزن بلوک، میانگین اندازه تراکنش، و فاصله زمانی تولید بلوک. اگر اندازه متوسط هر تراکنش را حدود چند صد بایت و فاصله بلوک را ۶۰۰ ثانیه در نظر بگیریم، ظرفیت نظری لایه پایه در حد چند تراکنش بر ثانیه\LTRfootnote{Transactions Per Second (TPS)} باقی می‌ماند. ارتقای شاهد جداگانه\LTRfootnote{Segregated Witness (SegWit)} با معرفی مفهوم وزن بلوک، ظرفیت مؤثر را بهتر کرد، اما ماهیت تک‌زنجیره‌ای و تکرار کامل اعتبارسنجی در همه نودها همچنان پابرجاست. بنابراین، حتی بهبودهای لایه اول نیز معمولاً با سقف‌های عملی مواجه می‌شوند.

از سوی دیگر، کاربردهای امروزی بلاکچین بسیار فراتر از انتقال ساده ارزش هستند. پرداخت خرد، تسویه بین صرافی‌ها، قراردادهای هوشمند، سیستم‌های رأی‌گیری، زنجیره تأمین، هویت دیجیتال، دارایی‌های توکنیزه‌شده و برنامه‌های مالی غیرمتمرکز نیازمند ظرفیت بالاتر، تأخیر کمتر و تجربه کاربری بهتر هستند. مقایسه کمی، عمق فاصله را نشان می‌دهد: شبکه‌های کارت پرداخت در بار معمول چند هزار تراکنش بر ثانیه پردازش می‌کنند و ظرفیت اسمی ده‌ها هزار تراکنش بر ثانیه دارند، در حالی که لایه پایه بیتکوین در حدود هفت تراکنش بر ثانیه و اتریوم در حدود چند ده تراکنش بر ثانیه باقی می‌مانند؛ یعنی فاصله‌ای بین دو تا چهار مرتبه بزرگی. در چنین شرایطی، اگر کاربر برای تأیید اولیه یک پرداخت خرد چند ده دقیقه منتظر بماند یا کارمزد او از ارزش خود تراکنش بیشتر شود، کاربردپذیری سامانه از بین می‌رود.

راهکارهای متعددی برای حل مسئله ارائه شده‌اند: افزایش اندازه بلوک، کاهش فاصله بلوک، شبکه لایتنینگ\LTRfootnote{Lightning Network}، زنجیره‌های جانبی، شاردینگ\LTRfootnote{Sharding}، رول‌آپ‌ها\LTRfootnote{Rollup} و تغییر الگوریتم اجماع. هر یک از این رویکردها بخشی از مسئله را حل می‌کنند، اما هزینه‌ها و فرض‌های خاص خود را دارند. افزایش اندازه بلوک ساده است، اما هزینه اجرای نود کامل را افزایش می‌دهد. شبکه لایتنینگ برای پرداخت‌های تکرارشونده و خرد بسیار مفید است، اما نیازمند مدیریت کانال و نقدینگی است. شاردینگ ظرفیت را افزایش می‌دهد، اما تراکنش‌های بین‌شاردی و امنیت هر شارد را پیچیده می‌کند. رول‌آپ‌ها در اکوسیستم اتریوم نقش مهمی یافته‌اند، اما به زیرساخت اثبات، داده‌دسترسی و اغلب مرتب‌ساز\LTRfootnote{Sequencer} متمرکز وابسته‌اند. زنجیره‌های جانبی انعطاف‌پذیرند، اما باید مسئله اتصال، پل، امنیت و اعتماد را حل کنند.

با این مقدمات، مسئله این پایان‌نامه را می‌توان به صورت یک مسئله طراحی مقید صورت‌بندی کرد. هدف، یافتن معماری \(\mathcal{S}\) است که توان عملیاتی کل را بیشینه کند، در حالی که چهار قید برقرار بماند:

\begin{latin}
\[
\max_{\mathcal{S}} \; TPS(\mathcal{S}) \quad \text{s.t.} \quad
\begin{cases}
Sec(\mathcal{S}) \geq Sec(C_0) - \epsilon & \text{(security anchor)}\\
Cost_{node}(\mathcal{S}) = O(TPS/N) & \text{(node cost)}\\
A: U \rightarrow \{1,\dots,N\} \text{ deterministic, public} & \text{(assignment)}\\
p_{doublespend}(\mathcal{S}) \approx p_{doublespend}(C_0) & \text{(safety)}
\end{cases}
\]
\end{latin}

قید نخست می‌گوید قطعیت نهایی تراکنش‌ها باید همچنان به امنیت زنجیره اصلی \(C_0\) گره بخورد و افت امنیتی ناشی از معماری جدید محدود و قابل تحلیل باشد. قید دوم می‌گوید هزینه هر نود نباید با ظرفیت کل رشد کند؛ بلکه باید بتوان با دنبال کردن زیرمجموعه‌ای از زنجیره‌ها، با هزینه‌ای متناسب با سهم همان زیرمجموعه مشارکت کرد، وگرنه مقیاس‌پذیری به قیمت تمرکز به دست آمده است. قید سوم می‌گوید تخصیص تراکنش به واحدهای پردازشی باید بدون هماهنگ‌کننده مرکزی، توسط هر نودی به طور مستقل و قطعی قابل محاسبه باشد. قید چهارم می‌گوید تقسیم بار نباید سطح حمله جدیدی برای دوبارخرج باز کند. راهکارهای موجود هر کدام دست‌کم یکی از این قیدها را نقض می‌کنند یا ضعیف می‌کنند: افزایش اندازه بلوک قید دوم را، زنجیره جانبی فدرال قید نخست را، و شاردینگ مبتنی بر کمیته، بسته به اندازه کمیته، قید نخست و چهارم را تحت فشار قرار می‌دهد.

مسئله این پایان‌نامه در همین نقطه شکل می‌گیرد: طراحی معماری‌ای که با حفظ یک زنجیره اصلی امن، پردازش تراکنش را به چند زنجیره موازی تقسیم کند؛ به‌گونه‌ای که هر زنجیره جانبی فقط بخشی از تراکنش‌ها را پردازش کند، اما نتیجه پردازش از طریق ارجاع رمزنگاری‌شده به زنجیره اصلی متصل شود. در این طرح، زنجیره اصلی همچنان محل امنیت نهایی و توافق جهانی است، ولی همه تراکنش‌ها مستقیماً در آن نوشته نمی‌شوند. کلید فنی طرح، انتخاب مدل خروجی‌های خرج‌نشده به عنوان مبنای تقسیم بار است: چون هر تراکنش فقط به خروجی‌های مشخصی وابسته است، افراز فضای خروجی‌ها به \(N\) قلمرو، تراکنش‌ها را به \(N\) کلاس تقریباً مستقل تقسیم می‌کند و بررسی دوبارخرج هر خروجی به یک قلمرو واحد محدود می‌شود. این جداسازی نقش‌ها می‌تواند گلوگاه تک‌زنجیره‌ای را کاهش دهد و در فصل سوم نشان داده می‌شود که تحت فرض‌های مشخص، هر چهار قید بالا را برآورده می‌کند.

\section{اهمیت و ضرورت پژوهش}

اهمیت این پژوهش از چند جنبه قابل بررسی است. نخست، از دید نظری، مسئله مقیاس‌پذیری بلاکچین یکی از نمونه‌های مهم تعارض میان امنیت، کارایی و عدم تمرکز است. طراحی هر راهکار مقیاس‌پذیر باید به این پرسش پاسخ دهد که کدام بخش از بار پردازشی کاهش می‌یابد، کدام فرض امنیتی اضافه می‌شود، و کدام بازیگر جدید ممکن است به نقطه تمرکز تبدیل شود. بنابراین، حتی اگر یک معماری به عدد تراکنش بر ثانیه بالایی برسد، بدون تحلیل امنیت و مدل اعتماد، نمی‌توان آن را راهکار کامل دانست.

دوم، از دید کاربردی، پذیرش گسترده بلاکچین بدون ظرفیت کافی امکان‌پذیر نیست. سامانه‌های پرداخت سنتی، پایگاه‌های داده بانکی، شبکه‌های کارت و زیرساخت‌های مالی متمرکز ظرفیت تراکنشی بالاتری نسبت به بسیاری از بلاکچین‌های عمومی دارند. اگر بلاکچین قرار باشد در کاربردهای عمومی و پرترافیک مورد استفاده قرار گیرد، باید بتواند حداقل بخشی از فاصله عملکردی با این سامانه‌ها را جبران کند.

سوم، از دید مهندسی نرم‌افزار، پیاده‌سازی یک معماری چندزنجیره‌ای فرصت مناسبی برای بررسی عملی محدودیت‌های تئوری است. بسیاری از طرح‌های مقیاس‌پذیری در سطح فرمول، افزایش خطی ظرفیت را نشان می‌دهند؛ اما در پیاده‌سازی واقعی، سربار هماهنگی، مدیریت حافظه، رقابت میان نخ‌های پردازشی\LTRfootnote{Thread}، انتشار شبکه‌ای بلوک‌ها و هزینه اعتبارسنجی امضاها باعث می‌شود ظرفیت عملی کمتر از مقدار نظری باشد. این پایان‌نامه تلاش می‌کند میان مدل نظری و بنچمارک\LTRfootnote{Benchmark} عملی تفکیک قائل شود.

\section{چالش سه‌گانه مقیاس‌پذیری، امنیت و غیرمتمرکز بودن}

در ادبیات بلاکچین، معمولاً از اصطلاح «سه‌گانه بلاکچین»\LTRfootnote{Blockchain Trilemma} برای توصیف تعارض میان سه ویژگی استفاده می‌شود: مقیاس‌پذیری، امنیت و غیرمتمرکز بودن. مقصود از مقیاس‌پذیری، توان شبکه برای افزایش ظرفیت با رشد کاربران و تراکنش‌هاست. امنیت به مقاومت سامانه در برابر دوبارخرج، بازنویسی تاریخچه، سانسور و حملات اجماع اشاره دارد. غیرمتمرکز بودن نیز به امکان مشارکت تعداد زیادی نود مستقل در اعتبارسنجی و تصمیم‌گیری مربوط است.

در بیتکوین، امنیت و غیرمتمرکز بودن بر ظرفیت ترجیح داده شده‌اند. فاصله بلوک طولانی و محدودیت اندازه بلوک باعث می‌شود نودهای بیشتری بتوانند شبکه را دنبال کنند و احتمال فورک‌های مکرر کاهش یابد. اما همین انتخاب‌ها ظرفیت را محدود می‌کنند. در مقابل، برخی زنجیره‌ها با کاهش شدید فاصله بلوک یا افزایش اندازه بلوک، توان عملیاتی بیشتری ارائه می‌دهند؛ ولی ممکن است نیاز سخت‌افزاری نودها افزایش یابد یا تمرکز تولید بلوک بیشتر شود. هدف معماری پیشنهادی این پایان‌نامه آن نیست که این سه‌گانه را به طور کامل حذف کند، بلکه می‌کوشد با تقسیم نقش‌ها، فشار تراکنشی را از زنجیره اصلی کم کند و همچنان از امنیت زنجیره اصلی به عنوان لنگر نهایی استفاده نماید.

\section{اهداف پژوهش}

هدف اصلی این پایان‌نامه، طراحی و ارزیابی یک معماری اجماع و پردازش تراکنش مقیاس‌پذیر برای بلاکچین‌های مبتنی بر مدل خروجی‌های خرج‌نشده است. اهداف ویژه عبارت‌اند از:

\begin{enumerate}
\item ارائه معماری چندزنجیره‌ای مبتنی بر زنجیره اصلی و چند زنجیره جانبی موازی.
\item تعریف تابع تخصیص قطعی برای نگاشت تراکنش‌ها یا خروجی‌های خرج‌نشده به زنجیره‌های جانبی.
\item حفظ نقش امنیتی زنجیره اصلی از طریق ثبت ارجاع بلوک‌های جانبی در آن.
\item کاهش تأخیر تأیید اولیه از طریق تولید بلوک سریع‌تر در زنجیره‌های جانبی.
\item تحلیل تئوری مقیاس‌پذیری افقی و رابطه توان عملیاتی با تعداد زنجیره‌های جانبی.
\item پیاده‌سازی نمونه اولیه در زبان راست\LTRfootnote{Rust} و بررسی ظرفیت محاسباتی آن.
\item شناسایی محدودیت‌ها، ریسک‌های امنیتی و مسیرهای توسعه آینده.
\end{enumerate}

\section{پرسش‌های پژوهش}

پرسش‌های اصلی پژوهش به صورت زیر تعریف می‌شوند:

\begin{itemize}
\item چگونه می‌توان پردازش تراکنش‌ها را بین چند زنجیره جانبی تقسیم کرد، بدون آنکه خطر دوبارخرج به شکل غیرقابل کنترل افزایش یابد؟
\item نقش زنجیره اصلی در معماری چندزنجیره‌ای پیشنهادی چیست و چگونه می‌تواند به عنوان لنگر امنیتی عمل کند؟
\item رابطه تعداد زنجیره‌های جانبی با توان عملیاتی کل سیستم چگونه است؟
\item چه تفاوتی میان توان عملیاتی نظری شبکه و توان عملیاتی بنچمارک محاسباتی پیاده‌سازی وجود دارد؟
\item مهم‌ترین ریسک‌های امنیتی، مهندسی و عملیاتی معماری پیشنهادی کدام‌اند؟
\end{itemize}

\section{نوآوری و سهم پژوهش}

سهم پژوهش در چند محور قابل بیان است. نخست، یک چارچوب مفهومی و رسمی برای ترکیب زنجیره اصلی اثبات کار با چند زنجیره جانبی موازی ارائه می‌شود. دوم، از تابع تخصیص مبتنی بر هش خروجی‌های خرج‌نشده برای توزیع تراکنش‌ها استفاده می‌شود تا هر خروجی به یک قلمرو اعتبارسنجی مشخص تعلق گیرد. سوم، مدل ریاضی توان عملیاتی ارائه می‌شود که نشان می‌دهد تحت فرض‌های تقارنی، ظرفیت کل با تعداد زنجیره‌ها رشد خطی دارد. چهارم، پیاده‌سازی نمونه اولیه در زبان راست انجام شده و نتایج آن به صورت تفکیک‌شده از مدل نظری تحلیل می‌شود. پنجم، برخلاف نگارش‌های کوتاه که صرفاً عدد توان عملیاتی را گزارش می‌کنند، این پایان‌نامه درباره تهدیدهای اعتبار، محدودیت‌های بنچمارک، سربار هماهنگی و ملاحظات امنیتی بحث می‌کند.

\section{دامنه و فرضیات پژوهش}

این پژوهش بر بلاکچین‌های عمومی یا نیمه‌عمومی مبتنی بر مدل خروجی‌های خرج‌نشده تمرکز دارد. فرض می‌شود تراکنش‌ها دارای ورودی‌ها و خروجی‌های قابل شناسایی هستند و هر خروجی خرج‌نشده می‌تواند با هش شناسه خود به یک زنجیره جانبی تخصیص یابد. همچنین فرض می‌شود زنجیره اصلی از اثبات کار استفاده می‌کند و قواعد انتخاب زنجیره بر اساس بیشترین کار انباشته است. پیاده‌سازی نمونه اولیه با هدف ارزیابی معماری و ظرفیت محاسباتی انجام شده و جایگزین یک شبکه عمومی تولیدی نیست. بنابراین، نتایج بنچمارک نباید بدون در نظر گرفتن محدودیت‌های شبکه، اندازه بلوک، انتشار بلوک، اعتبارسنجی امضا و رفتار مهاجم به عنوان توان واقعی شبکه تلقی شوند.

\section{بسط مسئله از منظر صف و گلوگاه}

برای فهم دقیق‌تر مسئله مقیاس‌پذیری، می‌توان زنجیره اصلی را مانند یک صف با ظرفیت خروجی محدود در نظر گرفت. تراکنش‌ها از سوی کاربران وارد استخر تراکنش می‌شوند. ماینرها در هر بازه زمانی یک بلوک با ظرفیت محدود می‌سازند و فقط بخشی از تراکنش‌ها را از صف خارج می‌کنند. اگر نرخ ورود تراکنش‌ها از نرخ خروج بیشتر شود، طول صف افزایش می‌یابد. افزایش طول صف دو پیامد دارد: نخست، تأخیر تأیید تراکنش‌ها بیشتر می‌شود؛ دوم، کاربران برای اولویت گرفتن کارمزد بالاتری پیشنهاد می‌کنند. این پدیده در دوره‌های ازدحام شبکه بیتکوین و اتریوم بارها مشاهده شده است.

از دید نظری صف، اگر \(\lambda\) نرخ ورود تراکنش و \(\mu\) نرخ پردازش زنجیره باشد، زمانی که \(\lambda < \mu\) سیستم پایدار است، اما هرچه \(\lambda\) به \(\mu\) نزدیک‌تر شود، تأخیر متوسط به شکل غیرخطی افزایش می‌یابد. اگر \(\lambda > \mu\) شود، صف بدون کران رشد می‌کند مگر آنکه کاربران از ارسال تراکنش منصرف شوند یا کارمزدها آن‌قدر بالا رود که تقاضا کاهش یابد. بنابراین، مقیاس‌پذیری فقط به معنای داشتن توان عملیاتی بالاتر در شرایط آزمایشگاهی نیست؛ بلکه به معنای نگه داشتن سیستم در محدوده پایدار تحت بار واقعی است.

معماری پیشنهادی این پایان‌نامه، از دید صف، یک صف مرکزی را به چند صف موازی تبدیل می‌کند. هر زنجیره جانبی استخر تراکنش مستقل دارد و بخشی از تراکنش‌ها را پردازش می‌کند. اگر توزیع بار میان زنجیره‌ها متعادل باشد، نرخ خروج کل تقریباً برابر مجموع نرخ خروج زنجیره‌های جانبی خواهد بود. این همان مبنای رشد خطی در مدل تئوری است. اما اگر بار نامتعادل شود یا تراکنش‌های بین‌زنجیره‌ای زیاد باشند، بخشی از این مزیت از بین می‌رود. به همین دلیل، در فصل سوم علاوه بر فرمول توان عملیاتی، بحث تعادل بار و تراکنش‌های بین‌زنجیره‌ای نیز اضافه شده است.

\section{تمایز میان مقیاس‌پذیری عمودی و افقی}

در سامانه‌های رایانه‌ای، مقیاس‌پذیری عمودی به معنای افزایش توان یک ماشین یا یک زنجیره است؛ برای مثال افزایش اندازه بلوک، استفاده از سخت‌افزار قوی‌تر یا بهینه‌سازی کد. مقیاس‌پذیری افقی\LTRfootnote{Horizontal Scalability} به معنای اضافه کردن واحدهای پردازشی مستقل است؛ مانند افزودن شارد\LTRfootnote{Shard}، زنجیره جانبی یا نود پردازشی. در بلاکچین‌های عمومی، مقیاس‌پذیری عمودی محدودیت جدی دارد، زیرا اگر نیاز سخت‌افزاری نود کامل بیش از حد بالا رود، تعداد مشارکت‌کنندگان مستقل کاهش می‌یابد و سامانه به سمت تمرکز حرکت می‌کند.

پیشنهاد این پایان‌نامه در اصل یک راهکار افقی است. به جای اینکه یک زنجیره را بزرگ‌تر و سنگین‌تر کنیم، چند زنجیره جانبی موازی ایجاد می‌کنیم. البته مقیاس‌پذیری افقی نیز رایگان نیست. باید مشخص شود هر زنجیره چه کسی را اعتبارسنج دارد، چگونه با زنجیره اصلی هماهنگ می‌شود، چگونه داده‌های خود را منتشر می‌کند و چگونه از حمله جلوگیری می‌شود. مزیت رویکرد افقی آن است که ظرفیت می‌تواند به صورت مرحله‌ای افزایش یابد و لازم نیست همه نودها در هر لحظه تمام تراکنش‌های همه زنجیره‌ها را پردازش کنند.

\section{تعریف دقیق‌تر مسئله پژوهش}

با توجه به مباحث بالا، مسئله پژوهش را می‌توان به صورت دقیق‌تر چنین بیان کرد: «طراحی معماری بلاکچینی مبتنی بر مدل خروجی‌های خرج‌نشده که بتواند با استفاده از مجموعه‌ای از زنجیره‌های جانبی موازی، نرخ پردازش تراکنش را نسبت به معماری تک‌زنجیره‌ای افزایش دهد، در حالی که امنیت نهایی تراکنش‌ها از طریق زنجیره اصلی اثبات کار تأمین شود و قواعد تخصیص تراکنش‌ها به زنجیره‌های جانبی به صورت قطعی، عمومی و قابل اعتبارسنجی باشد.»

این تعریف چند قید مهم دارد. نخست، معماری باید با مدل خروجی‌های خرج‌نشده سازگار\LTRfootnote{UTXO-compatible} باشد. دوم، امنیت نهایی باید به زنجیره اصلی متصل باشد. سوم، تخصیص تراکنش نباید به تصمیم متمرکز وابسته باشد. چهارم، افزایش توان عملیاتی باید با تحلیل محدودیت‌ها همراه باشد. پنجم، پیاده‌سازی باید آن‌قدر دقیق باشد که امکان اندازه‌گیری، بازتولید و نقد فنی فراهم شود.

\section{ساختار پایان‌نامه}

فصل دوم به مبانی نظری و مرور مطالعات انجام‌شده اختصاص دارد. در این فصل، بیتکوین، اثبات کار، مدل خروجی‌های خرج‌نشده، لایه دوم، شاردینگ، زنجیره‌های جانبی، رول‌آپ‌ها و روش‌های افزایش ظرفیت بررسی می‌شوند. فصل سوم معماری پیشنهادی را به صورت رسمی و فنی معرفی می‌کند و شامل مدل سیستم، ساختار داده‌ها، الگوریتم‌ها، تحلیل امنیت و مدل مقیاس‌پذیری است. فصل چهارم نتایج پیاده‌سازی و ارزیابی را ارائه می‌دهد و تفاوت میان مدل نظری و بنچمارک عملی را توضیح می‌دهد. فصل پنجم جمع‌بندی، محدودیت‌ها و پیشنهادهای آینده را بیان می‌کند.
